% =================================================
% Домашняя работа ученика
% =================================================

\packtitle{Домашняя работа}{Серединные перпендикуляры и скрытая окружность}
\studentnameblock

\begin{routebox}[Что нужно уметь]
Перед выполнением вспомните четыре связки:
\[
\begin{aligned}
&\text{середина}+\text{перпендикуляр}
&&\Longrightarrow \text{серединный перпендикуляр},\\
&\text{два серединных перпендикуляра}
&&\Longrightarrow \text{центр описанной окружности},\\
&\text{центр}+\text{перпендикуляр к хорде}
&&\Longrightarrow \text{деление хорды пополам},\\
&\text{вписанный угол}
&&\Longrightarrow \text{центральный угол вдвое больше}.
\end{aligned}
\]
\end{routebox}



\section{Базовые конструкции}

\begin{homeworktask}[Серединный перпендикуляр]{1}
Точка \(K\) — середина отрезка \(PQ\). Через \(K\) проведена прямая,
перпендикулярная \(PQ\), и на ней взята точка \(O\). Докажите, что
\[
OP=OQ.
\]
\noindent
\begin{minipage}[t]{0.43\linewidth}
\drawingbox[36mm]
\end{minipage}\hfill
\begin{minipage}[t]{0.53\linewidth}
\answerline[2]
\end{minipage}
\end{homeworktask}

\Needspace{15\baselineskip}
\begin{homeworktask}[Центр описанной окружности]{2}
В треугольнике \(XYZ\) серединные перпендикуляры к сторонам \(XY\) и
\(XZ\) пересекаются в точке \(O\).

\begin{enumerate}
  \item Докажите, что \(OY=OZ\).
  \item Как называется точка \(O\)?
\end{enumerate}
\noindent
\begin{minipage}[t]{0.43\linewidth}
\drawingbox[36mm]
\end{minipage}\hfill
\begin{minipage}[t]{0.53\linewidth}
\answerline[3]
\end{minipage}
\end{homeworktask}

\begin{homeworktask}[Центральный и вписанный углы]{3}
Точки \(A,B,C\) лежат на окружности с центром \(O\). Известно, что
\[
\angle ACB=57^\circ.
\]
Найдите \(\angle AOB\) и углы при основании равнобедренного треугольника
\(AOB\).
\noindent
\begin{minipage}[t]{0.43\linewidth}
\drawingbox[36mm]
\end{minipage}\hfill
\begin{minipage}[t]{0.53\linewidth}
\answerline[3]
\end{minipage}
\end{homeworktask}

\newpage
\begin{homeworktask}[Расстояние от центра до хорды]{4}
В окружности с центром \(O\) проведена хорда \(BC\). Перпендикуляр
\(OE\) опущен на \(BC\). Известно, что
\[
BC=12\text{ см},
\qquad
OE=6\text{ см}.
\]
Найдите \(\angle OBC\).

\begin{hintbox}
Сначала объясните, почему \(BE=EC\).
\end{hintbox}
\noindent
\begin{minipage}[t]{0.43\linewidth}
\drawingbox[35mm]
\end{minipage}\hfill
\begin{minipage}[t]{0.53\linewidth}
\answerline[3]
\end{minipage}
\end{homeworktask}

\section{Комбинируем идеи}

\begin{homeworktask}[Два угла из разных конструкций]{5}
В треугольнике \(BNC\) точка \(M\) — центр описанной окружности.
Известно, что
\[
\angle BCN=62^\circ,
\qquad
\operatorname{\rho}(M,BC)=\frac12BC.
\]
Луч \(BM\) проходит внутри угла \(NBC\). Найдите \(\angle NBC\) (\(\rho(M,BC)\)- расстояние от точки \(M\)до прямой \(BC\)).


\noindent
\begin{minipage}[t]{0.43\linewidth}
\drawingbox[42mm]
\end{minipage}\hfill
\begin{minipage}[t]{0.53\linewidth}
\answerline[4]
\end{minipage}
\end{homeworktask}

\Needspace{24\baselineskip}
\begin{homeworktask}[Аналог основной задачи]{6}
В трапеции \(ABCD\) выполнено
\[
AD\parallel BC,
\qquad
BC=2AD.
\]
Внутри трапеции взята точка \(M\) так, что
\[
\angle BAM=\angle CDM=90^\circ.
\]
Расстояние от точки \(M\) до прямой \(BC\) равно \(AD\), а
\[
\angle BCD=68^\circ.
\]
Найдите \(\angle ABC\).
\drawingbox[62mm]
\answerline[4]
\end{homeworktask}

\Needspace{22\baselineskip}
\section{Анализ и обобщение}

\begin{homeworktask}[Найдите ошибку]{7}
После достроения трапеции до треугольника \(BNC\) ученик опустил
\(ME\perp BC\) и написал:

\begin{quote}
«Так как через точку \(M\) к прямой \(BC\) можно провести только один
перпендикуляр, то точки \(N,M,E\) лежат на одной прямой».
\end{quote}

Объясните ошибку. Какое дополнительное утверждение потребовалось бы доказать,
чтобы рассуждение стало верным?
\answerline[3]
\end{homeworktask}

\begin{homeworktask}[Общая формула]{8}
В условиях задачи 6 заменим \(68^\circ\) на произвольный угол
\[
\angle BCD=\alpha.
\]
Выразите \(\angle ABC\) через \(\alpha\).

Заполните маршрут:
\[
\angle MBC=\shortanswer{18mm},
\qquad
\angle MBN=\shortanswer{28mm},
\]
\[
\angle ABC=\shortanswer{42mm}.
\]
\answerline[2]
\end{homeworktask}
